% ==============================================================================
% Section 2: Theoretical Framework
% AR-DeC: A UML-Driven Framework for Mobile AR-Enhanced Vocabulary Acquisition
% ==============================================================================

\section{Theoretical Framework}
\label{sec:theoretical}

This section establishes the theoretical foundations upon which the AR-DeC framework is constructed. We examine five intersecting domains---mobile augmented reality in education, vocabulary acquisition theory, intelligent tutoring systems, gamification, and UML-based pedagogical modeling---and synthesize them into a coherent rationale for our integrated design.

% ------------------------------------------------------------------------------
\subsection{Mobile Augmented Reality in Education: Foundations and Challenges}
\label{subsec:ar_education}

Augmented reality (AR) is broadly defined as a technology that superimposes computer-generated content onto the user's perception of the real world, satisfying three foundational criteria: the combination of real and virtual elements, real-time interactivity, and three-dimensional registration within the physical environment \cite{Azuma1997}. Unlike virtual reality, which immerses the user in an entirely synthetic space, AR preserves the primacy of the real-world context while enriching it with supplementary digital information. This characteristic makes AR particularly well-suited to educational settings, where learners benefit from maintaining situational awareness while receiving contextual support. In the mobile variant of AR, the smartphone or tablet camera serves as the sensory gateway, and the device's inertial measurement unit, GPS, and computer vision algorithms cooperate to anchor virtual overlays to physical surfaces in real time.

The pedagogical promise of mobile AR rests on several well-documented affordances. First, spatial anchoring allows digital annotations to be registered directly onto the objects or texts that learners are examining, thereby reducing the split-attention effect that arises when information sources are physically separated \cite{Wu2013}. Second, contextual overlay enables just-in-time delivery of explanations, definitions, and visualizations precisely at the point of need, supporting situated cognition theory. Third, the embodied and multisensory nature of AR interactions---pointing a camera, moving around an object, tapping on overlaid elements---has been shown to promote deeper encoding through motoric engagement \cite{Akcayir2017}. Meta-analytic syntheses have confirmed these benefits quantitatively: Garz\'{o}n and Acevedo \cite{Garzon2019} reported moderate-to-large effect sizes across 64 comparative studies, while Ib\'{a}\~{n}ez and Delgado-Kloos \cite{Ibanez2018} found that AR-enhanced learning environments consistently outperformed conventional instruction on both knowledge gain and motivational measures.

Despite this encouraging evidence, the literature also reveals persistent challenges that limit the educational impact of current mobile AR implementations. Cognitive overload can emerge when excessive or poorly designed overlays compete for attentional resources, paradoxically undermining the learning gains that AR is intended to produce \cite{Bacca2014}. Device limitations---including restricted field of view, battery drain, thermal throttling, and inconsistent tracking accuracy across hardware---further constrain the reliability of mobile AR experiences in authentic classroom conditions \cite{Dey2018}. Perhaps most critically, a recurrent finding across systematic reviews is that the majority of educational AR applications function primarily as content-delivery vehicles: they overlay static information onto the environment without adapting to individual learner needs, tracking mastery progression, or scaffolding cognitive engagement in a principled manner. This pedagogical shallowness reflects a disconnect between the technological capability of AR and the learning-science principles that should govern its deployment. There is therefore a compelling need for AR systems that are grounded in established learning theories, integrate adaptive mechanisms responsive to individual learner states, and embed motivational strategies that sustain engagement beyond the novelty effect. The AR-DeC framework addresses this need by coupling mobile AR with intelligent tutoring and gamification within a formally modeled pedagogical architecture.

% ------------------------------------------------------------------------------
\subsection{Vocabulary Acquisition and Reading Comprehension}
\label{subsec:vocabulary}

Vocabulary knowledge is widely recognized as the single strongest predictor of reading comprehension across first- and second-language contexts \cite{Nation2013, Stahl2005}. The relationship is reciprocal and cumulative: readers with larger vocabularies comprehend more text, and comprehending more text in turn exposes them to a greater range of new words, creating a virtuous cycle that Nation terms the ``Matthew effect'' in vocabulary growth. Conversely, students with limited vocabularies encounter a cascade of comprehension failures that compound over time, making early and effective vocabulary intervention a pedagogical priority of the highest order.

Research on vocabulary acquisition distinguishes between two fundamental modes of word learning. Incidental acquisition occurs when learners infer word meanings as a by-product of engaging with text for communicative purposes, while intentional learning involves the deliberate study of word forms, meanings, and usage patterns through explicit instruction or self-directed activities \cite{Hulstijn2001}. Both pathways contribute to vocabulary growth, but their effectiveness depends critically on the depth of cognitive processing that the learner invests. The involvement load hypothesis advanced by Laufer and Hulstijn \cite{Laufer2010} posits that vocabulary retention is a function of three processing components---need, search, and evaluation---and that tasks eliciting higher involvement loads produce more durable word knowledge. This framework aligns with the broader levels-of-processing tradition, which holds that deeper semantic engagement with a word (e.g., generating a sentence, analyzing morphological structure, or comparing the word with near-synonyms) yields stronger memory traces than shallow rehearsal of form-meaning pairs.

The advent of ubiquitous mobile devices has given rise to Mobile-Assisted Language Learning (MALL), a paradigm in which smartphones and tablets serve as platforms for vocabulary practice through flashcard applications, spaced-repetition systems, and multimedia dictionaries \cite{Kukulska2009, Chee2017}. MALL tools offer compelling advantages in terms of portability, immediacy, and personalization; however, the majority of existing applications present vocabulary in decontextualized lists divorced from the authentic texts in which learners encounter unfamiliar words. Empirical evidence consistently demonstrates that contextual vocabulary learning---encountering and processing words within meaningful discourse---produces superior retention and transfer compared to isolated word-list study \cite{Webb2007, Schmitt2008}. This finding underscores a critical gap in the current landscape: what is needed is an in-situ, adaptive vocabulary support mechanism that assists learners at the precise moment they encounter an unknown word during authentic reading, provides contextually grounded explanations calibrated to their current proficiency, and scaffolds progressive deepening of word knowledge over time. The AR-DeC framework is designed to fill precisely this gap by using optical character recognition to capture text from physical materials and delivering adaptive, AR-overlaid vocabulary support in real time.

% ------------------------------------------------------------------------------
\subsection{Intelligent Tutoring Systems for Adaptive Learning}
\label{subsec:its}

Intelligent Tutoring Systems (ITS) are computer-based instructional environments that provide individualized feedback, hints, and content selection by modeling both the knowledge domain and the learner's evolving cognitive state \cite{VanLehn2011}. The canonical architecture of an ITS comprises four interacting components: the \textit{domain model}, which encodes the knowledge structure and expert-level competence within the subject area; the \textit{student model}, which maintains a dynamic representation of each learner's current knowledge, misconceptions, and mastery levels; the \textit{tutoring model}, which embodies pedagogical strategies for selecting appropriate instructional actions given the student's state; and the \textit{interface model}, which governs how content and feedback are presented to the learner. This four-component architecture has proven remarkably versatile, supporting ITS implementations across domains ranging from mathematics and physics to programming and medical diagnosis.

The effectiveness of ITS has been rigorously quantified through meta-analytic research. VanLehn \cite{VanLehn2011} reported that well-designed ITS produce an average effect size of $d = 0.76$ on learning outcomes---an improvement nearly equivalent to that of one-on-one human tutoring ($d = 0.79$) and substantially exceeding that of conventional classroom instruction. This near-parity with human tutors is attributable to the ITS's capacity for continuous, fine-grained adaptation: unlike static instructional materials that present the same content to all learners regardless of prior knowledge, an ITS continuously updates its estimate of the learner's mastery and selects the next instructional action accordingly.

A foundational technique for student modeling in ITS is Bayesian Knowledge Tracing (BKT), introduced by Corbett and Anderson \cite{Corbett1994}. BKT models each knowledge component as a binary latent variable (learned or unlearned) and updates the probability of mastery after each learner interaction using Bayesian inference. Specifically, the probability that a student has learned a knowledge component after the $n$-th opportunity is given by:

\begin{equation}
\label{eq:bkt}
P(L_n) = P(L_{n-1}) + \big(1 - P(L_{n-1})\big) \cdot P(T)
\end{equation}

\noindent where $P(L_{n-1})$ is the posterior probability of mastery after the previous opportunity, $P(T)$ is the probability of transitioning from the unlearned to the learned state on any given opportunity, and the posterior is further conditioned on whether the observed response was correct or incorrect via the guess ($P(G)$) and slip ($P(S)$) parameters. In the vocabulary domain, BKT can be applied at the per-word level, maintaining an individual mastery estimate for every word a student encounters and enabling the system to prioritize review of words that remain below a mastery threshold.

Complementing the student model, the tutoring model in AR-DeC draws on Bloom's revised taxonomy \cite{Bloom1956, Krathwohl2002} to scaffold vocabulary instruction across progressively deeper cognitive levels. At the \textit{Remember} level, the system presents basic definitions and translations; at the \textit{Understand} level, it provides contextual examples, paraphrases, and visual associations; at the \textit{Apply} level, it poses fill-in-the-blank exercises and multiple-choice quizzes that require productive use of the word; and at the \textit{Analyze} level, it prompts the learner to examine morphological structure, compare near-synonyms, or identify semantic relationships with previously learned vocabulary. The BKT mastery estimate determines which Bloom's level is appropriate for each word at each encounter, creating a tightly coupled adaptive loop between assessment and instruction. Recent advances in large language models (LLMs) have further expanded the explanatory capacity of ITS, enabling the generation of contextually rich, natural-language definitions, example sentences, and etymological notes that surpass the quality achievable by template-based approaches \cite{Chen2020, Lin2023}. AR-DeC leverages these capabilities to produce vocabulary explanations that are both pedagogically principled and linguistically authentic.

% ------------------------------------------------------------------------------
\subsection{Gamification in Educational Contexts}
\label{subsec:gamification}

Gamification is defined as the application of game design elements and mechanics in non-game contexts to enhance user engagement, motivation, and behavioral outcomes \cite{Deterding2011}. In educational settings, gamification typically employs a constellation of mechanics including points, badges, leaderboards, levels, streaks, quests, and progress indicators to create a reward structure that reinforces desired learning behaviors \cite{Hamari2014}. The appeal of gamification lies in its capacity to activate both extrinsic motivational pathways---through tangible rewards and social comparison---and, when carefully designed, intrinsic motivational pathways through autonomy, competence feedback, and meaningful choice.

The theoretical alignment between gamification and the ARCS motivational model proposed by Keller \cite{Keller2010} provides a principled basis for designing gamified learning experiences. The ARCS model identifies four conditions necessary for sustained motivation: \textit{Attention}, captured through perceptual arousal and inquiry stimulation; \textit{Relevance}, established by connecting learning activities to the learner's goals and experiences; \textit{Confidence}, fostered by providing achievable challenges with clear success criteria; and \textit{Satisfaction}, achieved through intrinsic reinforcement and rewarding outcomes. In the AR-DeC framework, these four conditions map directly onto the system's core features: AR visual overlays and animated feedback capture \textit{Attention}; scanning personally encountered vocabulary from the learner's own textbooks establishes \textit{Relevance}; adaptive difficulty calibrated by BKT mastery estimates builds \textit{Confidence} by ensuring that challenges are neither trivially easy nor frustratingly difficult; and the gamification layer---comprising points for scanning and reading, badges for milestone achievement, streak bonuses for consecutive daily engagement, and combo multipliers for sustained quiz accuracy---delivers \textit{Satisfaction} through immediate and cumulative reward signals.

Meta-analytic evidence supports the effectiveness of gamification in educational contexts. Sailer and Homner \cite{Sailer2017} found significant positive effects of gamification on both cognitive learning outcomes and motivational variables, with the strongest effects observed when multiple mechanics were combined and when the game elements were meaningfully integrated with the learning task rather than superficially appended. Hamari, Koivisto, and Sarsa \cite{Hamari2014} similarly reported predominantly positive results, though they cautioned that effects were moderated by the quality of implementation and the specific population studied. A recognized risk in gamification design is the potential for over-gamification, in which an excess of extrinsic rewards can crowd out intrinsic motivation, reduce learners' sense of autonomy, and orient attention toward reward accumulation rather than learning. AR-DeC mitigates this risk through several design choices: points are awarded for cognitively meaningful actions (scanning, reading explanations, answering quiz questions) rather than trivial interactions; badge criteria require demonstrated mastery rather than mere activity volume; and streak mechanics with daily caps encourage regular but not compulsive engagement, sustaining long-term participation without fostering unhealthy usage patterns.

% ------------------------------------------------------------------------------
\subsection{Pedagogical Scenario Modeling with UML}
\label{subsec:uml}

The design of technology-enhanced learning environments that integrate multiple pedagogical strategies---adaptive tutoring, gamified engagement, and augmented reality presentation---demands a modeling formalism capable of capturing both the structural and behavioral complexity of the system while remaining accessible to the diverse stakeholders involved in educational technology development. Informal descriptions, however detailed, are prone to ambiguity and incompleteness; they fail to provide the conceptual rigor necessary for verifying that a system's design faithfully instantiates the intended pedagogical principles. There is consequently a need for formal modeling approaches that can bridge the gap between pedagogical theory and software implementation, providing a shared vocabulary through which instructional designers, learning scientists, and software engineers can collaborate effectively.

The Unified Modeling Language (UML) has emerged as the de facto standard notation for specifying, visualizing, and documenting the architecture of software systems \cite{Booch2005, Rumbaugh2004}. Originally developed for general-purpose software engineering, UML offers a rich set of diagram types that collectively capture both the static structure and dynamic behavior of complex systems. Its visual nature facilitates communication among heterogeneous teams, its standardized semantics ensure unambiguous interpretation, and its modular diagram taxonomy allows modelers to present different viewpoints of the same system at appropriate levels of abstraction. These properties make UML an attractive formalism for educational technology design, where the system under consideration spans technical infrastructure, pedagogical logic, and user interaction patterns.

In the AR-DeC framework, we employ four complementary UML diagram types to model the pedagogical scenario comprehensively. The \textit{class diagram} captures the static structural relationships among key entities---learner profiles, vocabulary items, mastery records, gamification elements, and AR presentation objects---defining the ontological backbone of the system. The \textit{use case diagram} delineates the functional capabilities exposed to each actor (student, system, teacher), establishing the boundaries of system responsibility and the entry points for user interaction. Together, these two static diagrams constitute the \textit{structural view} of the framework, specifying \textit{what} the system contains and \textit{who} interacts with it. The \textit{activity diagram} models the end-to-end pedagogical workflow---from text scanning through OCR processing, ITS-driven explanation generation, AR overlay rendering, gamification reward calculation, and student model update---capturing the control flow, decision points, and parallel processes that govern the learner's experience. The \textit{sequence diagram} details the temporal ordering of messages exchanged among system components during representative interaction scenarios, making explicit the coordination protocols that ensure pedagogical coherence. Together, these dynamic diagrams constitute the \textit{behavioral view}, specifying \textit{how} the system operates and \textit{when} each component contributes to the learning process.

This dual-view modeling strategy---structural and behavioral---follows the methodology established by Ouariach et al. in their UML-driven approach to collaborative e-learning system design \cite{Ouariach2023, Ouariach2025, Ouariach2026}. Their work demonstrated that formalizing pedagogical scenarios through UML diagrams yields several tangible benefits: improved stakeholder alignment during the design phase, enhanced traceability from pedagogical requirements to implementation artifacts, and greater reusability of design patterns across related educational technology projects. Hummel et al. \cite{Hummel2004} similarly argued for the adoption of educational modeling languages that combine pedagogical expressiveness with computational formalizability, while Retbi et al. \cite{Retbi2012} showed that visual modeling notations reduce design ambiguity and accelerate iterative refinement in educational technology development. Our contribution extends this line of work by adapting UML-based pedagogical scenario modeling to the specific demands of mobile augmented reality systems, which introduce additional complexity through real-time sensor integration, on-device inference, and spatially registered content delivery. With the theoretical foundations of a coherent pedagogical design now established, we shift our attention to formalizing the AR-DeC framework at a more concrete architectural level through the UML diagrams presented in the following section.
